\section{Kết luận}

Khoá luận này đã đề xuất và triển khai một kiến trúc GraphRAG cho bài toán khám phá và truy vấn tri thức trong miền nghệ thuật Chèo. Trên cơ sở phân tích hạn chế của các hệ thống Retrieval-Augmented Generation truyền thống - vốn dựa trên truy xuất văn bản theo tương đồng véc-tơ và khó đáp ứng các truy vấn đòi hỏi suy luận đa bước - nghiên cứu đã xây dựng một hệ thống tích hợp đồ thị tri thức vào quy trình truy xuất và tạo sinh, đồng thời đánh giá có hệ thống qua các thí nghiệm định lượng trên benchmark và kiểm nghiệm với người dùng thực. Các kết quả thực nghiệm xác nhận tính khả thi và hiệu quả của hướng tiếp cận này trong một miền tri thức có cấu trúc quan hệ phức tạp.

Đóng góp thứ nhất là việc xây dựng và chuẩn hoá một đồ thị tri thức chuyên biệt cho nghệ thuật Chèo. Từ ontology gốc được kế thừa, đồ thị được chuyển đổi sang mô hình property graph trên Neo4j, mở rộng với các phân lớp chi tiết về loại vai diễn và yếu tố biểu diễn, đồng thời được bổ sung 14 luật suy diễn quan hệ gián tiếp và ba mức chỉ mục để phục vụ các truy vấn chuyên sâu. Đồ thị tri thức cuối cùng gồm bảy vở chèo cổ, 38 nhân vật, 41 diễn viên, 15 trích đoạn và 27 phiên bản biểu diễn, tạo nền tảng dữ liệu có cấu trúc cho các nghiên cứu tiếp theo trong lĩnh vực số hoá di sản văn hoá truyền thống.

Đóng góp thứ hai là kiến trúc GraphRAG với hai phân hệ G-Retrieval và G-Generation được thiết kế tích hợp chặt chẽ. Phân hệ G-Retrieval triển khai truy vấn đa mức độ trên đồ thị - từ nút đơn lẻ, bộ ba quan hệ, đường dẫn đến tiểu đồ thị - đồng thời tổ chức ngữ cảnh theo bốn định dạng biểu diễn khác nhau để thích ứng với nhiều dạng câu hỏi. Phân hệ G-Generation áp dụng ba chiến lược Pre-Generation, Mid-Generation và Post-Generation, cân bằng giữa chất lượng câu trả lời và chi phí gọi mô hình ngôn ngữ. Kết quả thực nghiệm trên 100 ca của bộ dữ liệu CheoBench cho thấy GraphRAG đạt điểm tổng hợp $0,835$ so với $0,595$ của RAG truyền thống, với khoảng cách lớn nhất nằm ở khối Context (Context Recall $0,964$ so với $0,418$) và khối Generation (Faithfulness $0,986$ so với $0,734$) - khẳng định vai trò của cấu trúc đồ thị trong việc đảm bảo ngữ cảnh có liên kết và đầy đủ cho mô hình ngôn ngữ.

Đóng góp thứ ba là bộ dữ liệu đánh giá CheoBench gồm 100 câu hỏi được biên soạn thủ công và chia thành ba nhóm Local, Community và Global dựa trên phạm vi truy xuất và độ phức tạp suy luận. Bộ dữ liệu đi kèm hệ thống độ đo đa chiều gồm chín thành phần thuộc ba khối IR, Context và Generation, cho phép tách biệt chất lượng truy xuất, chất lượng ngữ cảnh và chất lượng câu trả lời cuối cùng trong quá trình phân tích hiệu năng. Bộ benchmark này không chỉ phục vụ đánh giá hệ thống được đề xuất trong nghiên cứu mà còn có thể tái sử dụng trong các nghiên cứu mở rộng về hệ thống hỏi đáp trên tri thức văn hoá tiếng Việt.

\section{Định hướng phát triển}
Mặc dù kiến trúc GraphRAG đề xuất đã chứng minh được tính khả thi và hiệu quả rõ rệt trong việc giảm thiểu hiện tượng ảo giác đối với miền tri thức nghệ thuật Chèo, nghiên cứu vẫn còn tồn tại một số hạn chế nhất định. Về mặt vận hành, hệ thống hiện có độ trễ trung vị khoảng 55 giây do chi phí tính toán lớn của cơ chế truy xuất đa mức độ, từ đó tạo ra trở ngại đối với các kịch bản yêu cầu tương tác liên tục hoặc phản hồi thời gian thực. Bên cạnh đó, quá trình đánh giá thực nghiệm vẫn còn giới hạn về mặt chuyên môn khi khảo sát mới được thực hiện trên 103 người dùng, trong đó hơn 74\% không có nền tảng chuyên sâu về nghệ thuật Chèo. Vì vậy, kết quả thu được hiện thiên về đánh giá trải nghiệm sử dụng và chất lượng diễn đạt hơn là khả năng kiểm chứng tri thức ở mức học thuật sâu. Ngoài ra, quy mô đồ thị tri thức hiện tại vẫn mang tính chất của một hệ thống thử nghiệm khái niệm khi mới được xây dựng trên phạm vi 7 vở chèo cổ tiêu biểu, do đó chưa phản ánh toàn diện tính đa dạng và độ phức tạp của kho tàng sân khấu dân gian Việt Nam.

Từ những hạn chế trên, các hướng phát triển trong tương lai sẽ tập trung vào việc tối ưu hiệu năng của hệ thống thông qua nghiên cứu cơ chế song song hóa cho bốn mức truy xuất và cải tiến bộ định tuyến truy vấn nhằm giảm thời gian phản hồi. Đồng thời, hệ thống cần được kiểm chứng chuyên môn sâu hơn với sự tham gia trực tiếp của các chuyên gia và nhà nghiên cứu trong lĩnh vực sân khấu truyền thống nhằm đánh giá chính xác hơn tính đúng đắn của tri thức và khả năng suy luận của mô hình. Xa hơn, nghiên cứu hướng tới việc mở rộng quy mô đồ thị tri thức và hoàn thiện kiến trúc GraphRAG phân tầng để ứng dụng cho các loại hình nghệ thuật truyền thống khác như Tuồng, Cải lương hay Dân ca, qua đó góp phần hỗ trợ công tác số hóa, bảo tồn và lan tỏa di sản văn hóa dân tộc bằng các công nghệ trí tuệ nhân tạo hiện đại.